\chapter{格式要求}
\label{chap:format-requirements}

学位论文每部分从新的一页开始，各部分要求如下：

\section{中文封面}
\label{sec:format-cover}

封面（cover）是学位论文的外表面，对论文起装潢和保护作用，并提供相关的信息。
我校申请博士、硕士和硕士专业学位的学位论文封面分别使用统一规定的不同封面。
封面中除已固定的内容外，其他需要填写的内容要求如下：

{\heiti 分类号：}暂空

{\heiti 密级：}按GB/T 7156《文献保密等级代码与标识》~\cite{SCSF00000734}标注。
公开论文可以标注“公开”，也可不标注；非公开论文标注“限制”、“秘密”或“机密”。
根据《南开大学关于研究生学位论文收藏和利用管理办法》（南发字〔2009〕23号文件）的规定，非公开论文须经申请、批准方能标注论文的密级（限制、秘密或机密），同时还应注明相应的保密年限（具体标注要求见\ref{sec:writing-classification}）。
非书面形式的涉密学位论文载体，要以能够明显识别的方式予以标注；学位论文电子文档中含有国家秘密内容的，应做出国家秘密标志，且标志与文档正文不可分离。

{\heiti 学校代码：}10055

{\heiti 学号：}填写本人学（申请）号。

{\heiti 论文题目（title，又称题名）：}应以简明词语恰当、准确地反映出论文最重要的特定内容，一般不宜超过25字，必要时可加论文副标题。

论文题目通常由名词性短语构成，应尽量避免使用不常用缩略词、首字母缩写字、字符、代号和公式等。

如论文题目内容层次很多，难以简化时，可采用论文题目和论文副标题相结合的方法，其中副标题起补充、阐明题目的作用。

示例1：斑马鱼和人的造血相关基因以及表观遗传学调控基因——进化、表达谱和功能研究

示例2：阿片镇痛的调控机制研究：Delta型阿片肽受体转运的调控机理及功能

{\heiti 培养单位：}指学位申请人所在学院（所）名称，应采用规范全称。
如哲学院、数学科学学院等。

{\heiti 一级学科：}学科名称以国务院学位委员会颁布的《研究生教育学科专业目录（2022年）》为准。

{\heiti 二级学科：}学科名称参照国务院学位委员会颁布的《授予博士、硕士学位和培养研究生的学科、专业目录》~\cite{REF00000002}或我校自主设置学科名称。

{\heiti 专业学位名称：}填写硕士专业学位名称。
如工商管理硕士、法律硕士等等。

{\heiti 论文作者：}填写作者姓名。

{\heiti 指导教师：}填写指导教师的姓名、职称（教授、研究员等）。

{\heiti 论文完成时间（年月）：}填写论文提交评审的时间。

\section{题名页}

题名页（title page）包含论文全部书目信息，单独成页（示例见附录A）。
主要内容规定如下：

{\heiti 中图分类号、UDC}

暂空

{\heiti 学校代码}

按照教育部批准的学校代码标注，应为“10055”。

{\heiti 密级}

标注同中文封面密级要求（具体标注见\ref{sec:writing-classification}）。

{\heiti 学位授予单位名称和学位论文类型}

保持“南开大学”题字不变。
根据学位论文类型填写“硕士学位论文”或“博士学位论文”。
请注意区分是学术学位还是专业学位（见附录A）。

{\heiti 题名}

题名要求同论文题目，应中英文对照。英文题名在中文题名下方。

题名在整篇学位论文中的不同地方出现时，应保持一致。

{\heiti 责任者}

责任者包括学生姓名，指导教师姓名、职称等。
校外导师应一并列出，由论文指导委员会指导论文写作的，应列出全部指导教师。
包括校内导师和校外导师。

如责任者姓名有必要附注汉语拼音时，遵照GB/T 16159《汉语拼音正词法基本规则》~\cite{SCSF00038519}著录。

{\heiti 申请学位}

包括申请的学位类别和级别，学位类别标注包括以下门类：哲学、经济学、法学、教育学、文学、历史学、理学、工学、农学、医学、军事学、管理学、艺术学。
学位级别标注包括硕士、博士。如哲学硕士、管理学博士。
专业学位直接标注其名称和层次，如工商管理硕士、临床医学博士等。

{\heiti 培养单位}

同\ref{sec:format-cover}中说明。

{\heiti 学科专业}

参照国务院学位委员会颁布的《研究生教育学科专业目录（2022年）》、《授予博士、硕士学位和培养研究生的学科、专业目录》~\cite{REF00000002}填写。
硕士专业学位填写领域名称，如无领域不必填此项。

{\heiti 研究方向}

指本学科专业范畴下的研究方向。

{\heiti 论文完成时间}

同\ref{sec:format-cover}中说明。

\section{勘误页（如有）}

学位论文如有勘误页，应另起页，放在题名页后。
在勘误页顶部应放置下列信息：题名、副题名（如有）、作者名。

\section{学位论文原创性声明和非公开学位论文标注说明}

本部分放在题名页之后另起页，内容见附录B，可直接将附录B复制到论文中，但要删除附录题目。
提交时学位论文原创性声明须有作者亲笔签名。
非公开学位论文标注说明须有南开大学学位评定委员会办公室盖章方为有效。

\section{学位论文使用授权书}

本部分放在学位论文原创性声明之后另起页，内容见附录C，可直接将附录C复制到论文中，但要删除附录题目。
提交时须有作者亲笔签名和导师亲笔签名。

学位论文应有使用授权说明。
学位申请人员应遵守学位授予单位有关保存、使用学位论文的相关规定。
学位授予单位有权保存学位论文的印刷本和电子版，并提供文献检索与阅览服务，以及向国家有关部门或机构送交学位论文的复印件和电子版。

\section{中文摘要}

学位论文应有摘要。
摘要应具有独立性和自含性，即不阅读论文的全文，就能获得必要的信息。
摘要的内容应包含与论文等同量的主要信息，供读者确定有无必要阅读全文，也可供二次文献采用。
摘要宜说明研究工作目的、方法、结果和结论等，重点是结果和结论。
摘要的撰写应符合GB/T 6447的规定。

以中文撰写的学位论文，中文摘要一般字数为300字~1000字，如遇特殊需要字数可略多。
摘要中宜尽量避免采用图、表、化学结构式、非公知公用的符号和术语。

每篇论文应选取3个~8个关键词，用显著的字符另起一行，排在摘要的下方。
关键词应体现论文特色，选用专业术语及词汇，在论文中有明确的出处。关键词的撰写应符合CY/T 173的规定。

为便于国际交流，宜有英文摘要，另起一页，置于中文摘要页的后面。
英文摘要的内容宜与中文摘要对应。
英文摘要的下方应标注与中文关健词对应的英文关键词。

以外文撰写的学位论文，应以同种语言撰写外文摘要。
除外文摘要外，应有中文摘要，置于外文摘要页的后面，并标注中文关键词。

\section{Abstract}

Abstract内容与中文摘要相对应。
一般不少于300个英文实词，篇幅以一页为宜。
如需要，字数可以略多。

\section{序言或前言（如有）}

学位论文的序言或前言一般是作者对本篇论文基本特征的简介，如说明研究工作的缘起、背景、主旨、目的、意义、编写体例，以及资助、支持、协作经过等。
这些内容也可以在正文引言（绪论）中说明。

\section{目录（目次）}

学位论文应有目录（目次）（table of contents）页，排在序言（或前言）之后，另起页。
目录是论文各章节标题的顺序列表，附有相应的起始页码。

\section{图和附表清单（如有）}

论文中如图表较多，可以分别列出清单置于目录页之后另起页。
图的清单应有序号、图题和页码。表的清单应有序号、表题和页码。

\section{符号、标志、缩略语等的注释表（如有）}

符号、标志、缩略语、首字母缩写、计量单位、自定义名词和术语等的注释说明，如需汇集，应编写成注释说明汇集表，可集中置于图表清单之后。
若上述符号使用数量不多，可以不设此部分，但必须在论文中出现时加以说明。

\section{正文}

正文是学位论文的主体部分，应从另页右页开始，每一章应另起页。
学位论文字数：建议博士学位论文10万字左右，硕士学位论文3万字左右。
在保证学位论文质量的前提下，各一级学科根据学科特点可自行规定论文字数。

正文部分应从另页右页开始，每一章应另起页。
正文部分通常包括引言（绪论）、主体和结论。

正文部分由于涉及的学科、选题、研究方法、结果表达方式等有很大的差异，不作统一的规定。
但是，应实事求是、客观真切、准备完备、合乎逻辑、层次分明、简练可读。
应该严格遵循学术规范，充分尊重和借鉴已有的学术成果，注重调查研究，在全面掌握相关研究资料和学术信息的基础上，精心设计研究方案，讲究科学方法。

\subsection{引言或绪言}

引言（绪论）应包括所研究领域的历史回顾和前人工作的综合述评，以及论文的研究目的、流程和方法等。

\subsection{具体章节}

本部分是论文作者的研究内容，是论文的核心。
各章之间互相关联，符合逻辑顺序。

\subsection{引文标注}

论文中引用的文献的标注方法遵照最新的GB/T 7714 《信息与文献参考文献著录规则》~\cite{SCSF00045714}，可采用顺序编码制，也可采用著者-出版年制，但全文必须统一。

\subsection{注释}

当论文中的字、词或短语，需要进一步加以说明，而又没有具体文献来源时，用注释（notes）。
注释一般在社会科学中用的较多。
注释的标注应符合CY/T 121《学术出版规范~注释》的规定。

注释可采用脚注、尾注或文中分散注释。
应控制论文中的注释数量，不宜过多。
由于学位论文篇幅较长，宜采用文中编号加脚注的方式，置于所注释正文所在页的底部。

\subsection{结论}

结论（最后一章）：是学位论文最终和总体的结论，应明确、精练、完整、准确，不是正文中各段的小结的简单重复。
论文的结论应包括论文的核心观点，着重阐述作者的创造性工作及所取得的研究成果在本学术领域的地位、作用和意义，交代研究工作的局限，提出未来工作的意见或建议。
如果不能导出一定的结论，也可以没有结论而进行必要的讨论。

\section{参考文献}

参考文献是文中引用的有具体文字来源的文献集合，其著录项目和著录格式应遵照GB/T 7714的规定执行。
参考文献应置于正文后，并另起页。
所有被引用文献均应列人参考文献中。

正文未被引用、仅拓展阅读或具有补充信息的文献，不应混入参考文献表，可集中置于附录，标题为“参考书目”。

参考文献表应置于正文后，并另起页。也可根据需要，在每页“脚注”中列出或在每章正文部分之后加入本章参考文献。

引文采用著者-出版年制标注时，参考文献表应按著者字顺和出版年排序。

\section{附录}

有些材料编入文章主体会有损于编排的条理性和逻辑性，或有碍于文章结构的紧凑和突出主题思想等，可将这些材料作为附录编排于全文的末尾。

附录放在正文之后另起页。
附录的序号用A，B，C，…系列，如附录A，附录B，…。
附录中的公式、图和表的编号分别用A1，A2，…系列；图A1，图A2，…系列；表A1，表A2，…系列。每个附录应有标题。

\section{索引（如有）}

如果需要，可以在附录后编排学位论文索引。
学位论文索引为学位论文编制的、由一系列款目组成，指向论文中的有检索意义的概念、语词及其创新内容等的信息检索工具，应依据 GB/T 41210《学位论文内容索引编制规则》的规定编制。

\section{致谢}

致谢是作者对该文章的形成作过贡献的组织或个人予以感谢的文字记载，语言要诚恳、恰当、简短。
致谢可单独成页，置于摘要页前，也可放在序言、前言或后记中。
包括国家科学基金，资助研究工作的奖学金基金、合同单位、资助或支持的企业、组织或个人；协助完成研究工作和提供便利条件的组织或个人；在研究工作中提出建议和提供帮助的人；给予转载和引用权的资料、图片、文献、研究和调查的所有者；其他应感谢的组织和个人。

\section{个人简历、在学期间发表的学术论文与研究成果}

包括作者的教育经历、攻读学位期间发表的学术论文及其他成果等。

\noindent 示例：

{\centering
YYYY-MM—YYYY-MM XX大学计算机系学士；

YYYY-MM—YYYY-MM XX大学攻读博士学位（直博）

获奖情况：

参加项目：

攻读博士学位期间发表的学术论文：

}
